\begin{examplebox}
	\textbf{\textcolor{CExample}{例 1（基础）}}

	\textbf{\textcolor{CExample}{题目：}}
	已知抛物线的焦点为$F(3,0)$，准线方程为$x=-3$，求抛物线方程。

	\textbf{\textcolor{CExample}{解题步骤：}}
	\begin{description}[style=nextline, leftmargin=2.8em, labelsep=0.5em, itemsep=0.45em]
		\item[\textbf{第一步（求参数$p$）：}] 由焦点在$x$轴正半轴，设标准方程$y^2=2px(p>0)$。焦点坐标为$(\frac{p}{2},0)$，准线$x=-\frac{p}{2}$。已知$(3,0)$，故$\frac{p}{2}=3$，得$p=6$。
			\par\textbf{理由：}
			标准方程焦点与准线形式。
		\item[\textbf{第二步（写方程）：}] 将$p=6$代入$y^2=2px$，得$y^2=12x$。
			\par\textbf{理由：}
			标准方程确定。
	\end{description}

	\textbf{\textcolor{CExample}{关键结论：}}
	\textbf{所求抛物线方程为$y^2=12x$。}

	\medskip
	\textbf{\textcolor{CExample}{例 2（稍变形）}}

	\textbf{\textcolor{CExample}{题目：}}
	设抛物线$y^2=8x$上一点$P$到焦点$F$的距离为$6$，求点$P$到准线的距离。

	\textbf{\textcolor{CExample}{解题步骤：}}
	\begin{description}[style=nextline, leftmargin=2.8em, labelsep=0.5em, itemsep=0.45em]
		\item[\textbf{第一步（识别参数）：}] 由$y^2=8x$得$2p=8$，$p=4$。焦点$F(2,0)$，准线$x=-2$。
			\par\textbf{理由：}
			标准方程系数对应。
		\item[\textbf{第二步（应用定义）：}] 根据抛物线第二定义，抛物线上点到焦点距离等于到准线距离。所以$P$到准线距离等于$|PF|=6$。
			\par\textbf{理由：}
			定义直接使用。
	\end{description}

	\textbf{\textcolor{CExample}{关键结论：}}
	\textbf{点$P$到准线的距离为$6$。}

	\medskip
	\textbf{\textcolor{CExample}{例 3（综合应用）}}

	\textbf{\textcolor{CExample}{题目：}}
	已知抛物线$C:y^2=2px(p>0)$的焦点为$F$，点$A(1,2)$在$C$上，求$p$和$|AF|$。

	\textbf{\textcolor{CExample}{解题步骤：}}
	\begin{description}[style=nextline, leftmargin=2.8em, labelsep=0.5em, itemsep=0.45em]
		\item[\textbf{第一步（代入求$p$）：}] 点$A(1,2)$在$C$上，代入$y^2=2px$得$4=2p\cdot1$，解得$p=2$。
			\par\textbf{理由：}
			点在曲线上满足方程。
		\item[\textbf{第二步（写标准方程）：}] 由$p=2$得抛物线方程为$y^2=4x$，焦点$F(1,0)$，准线$x=-1$。
			\par\textbf{理由：}
			标准方程对应关系。
		\item[\textbf{第三步（利用定义求距离）：}] 根据第二定义，$|AF|$等于点$A$到准线的距离。$A(1,2)$到$x=-1$的距离为$|1-(-1)|=2$，所以$|AF|=2$。
			\par\textbf{理由：}
			定义转化简化计算。
	\end{description}

	\textbf{\textcolor{CExample}{关键结论：}}
	\textbf{$p=2$，$|AF|=2$。}
\end{examplebox}
